% ============================================================================
% Fichier  : partie2/P02_C08_cobit_capacite_mise_en_oeuvre.tex
% Titre    : Mesurer et Mettre en Œuvre : Capacité, Maturité et Feuille de Route
% Auteur   : MINKA MI NGUIDJOÏ Thierry Emmanuel
% Version  : 1.0 -- octobre 2026
% ============================================================================
% Source : COBIT 2019 Framework: Introduction and Methodology, chapitres 6
% et 8 (dossier Normes), paraphrasés. Laboratoire associé : GCI-LAB-04
% (évaluation de capacité de DSS02 et BAI06).
% ============================================================================

\chapter{Mesurer et Mettre en Œuvre : Capacité, Maturité et Feuille de Route}
\label{chap:cobit-capacite}
\chaptermark{Capacité, maturité, mise en œuvre}

\epigraph{%
    « Ce qui ne se mesure pas ne s'améliore pas.\\
    Ce qui se mesure mal s'améliore dans la mauvaise direction. »%
}{Adage des praticiens de la qualité}

\niveauChapitre{Intermédiaire}{%
    Maîtrise du Chapitre~\ref{chap:cobit-conception} (facteurs de
    conception et niveaux de capacité cibles).%
}

\begin{boxobjectifs}
\begin{itemize}
    \item Appliquer le modèle de capacité des processus (niveaux 0 à 5).
    \item Distinguer capacité d'un processus et maturité d'un domaine.
    \item Évaluer un niveau de capacité à partir de pièces justificatives.
    \item Dérouler le cycle de mise en œuvre en sept phases.
    \item Construire une feuille de route d'amélioration.
\end{itemize}
\end{boxobjectifs}

\bigskip

% ============================================================================
\section*{Question directrice : où en sommes-nous vraiment~?}
% ============================================================================

Au Chapitre~\ref{chap:cobit-objectifs}, Mme~EKAMBI, directrice de l'audit
interne de la BANQUE SANAGA, avait présenté au conseil un tableau de
quarante cases « conforme / non conforme ». Au Chapitre~\ref{chap:cobit-conception},
le comité stratégique a fixé, pour dix objectifs prioritaires, une
capacité cible. Il reste à répondre à deux questions~: \emph{quelle est la
capacité réelle de chaque processus aujourd'hui~?} et \emph{comment combler
l'écart~?}

La première question relève de la \textbf{mesure de la performance}
(sections~\ref{sec:c08-principes} à~\ref{sec:c08-evaluer}). La seconde
relève de la \textbf{mise en œuvre} (sections~\ref{sec:c08-phases}
et~\ref{sec:c08-feuille}). Le laboratoire GCI-LAB-04 de NDAP SUKULU fait
évaluer la capacité de deux processus de la banque à partir d'un dossier
de pièces.

\separateur

% ============================================================================
\section{Mesurer la Performance dans COBIT}
\label{sec:c08-principes}
% ============================================================================

\subsection{Cinq principes}

COBIT~2019 désigne par \emph{gestion de la performance} l'ensemble des
méthodes qui disent dans quelle mesure le système de gouvernance et ses
composantes fonctionnent, et comment les améliorer. Il fixe cinq
principes à cette gestion~\cite{isaca2018cobitintro}~: elle doit être
\textbf{simple} à comprendre et à utiliser~; \textbf{cohérente} avec le
modèle conceptuel, et applicable à tous les types de composantes~;
\textbf{fiable}, c'est-à-dire produire des résultats reproductibles et
pertinents~; \textbf{flexible}, pour s'adapter aux priorités de chaque
organisation~; et capable de soutenir \textbf{différents types
d'évaluation}, de l'autoévaluation à l'audit formel.

\subsection{Un modèle aligné sur CMMI}

Le modèle de COBIT~2019 s'aligne largement sur la version~2.0 du CMMI et
l'étend. Il repose sur deux notions~:
\begin{itemize}
    \item la \textbf{capacité}, qui se mesure \emph{processus par
          processus}~: chaque activité d'un processus est rattachée à un
          niveau de capacité~;
    \item la \textbf{maturité}, qui se mesure \emph{domaine d'intérêt par
          domaine d'intérêt}~: elle est atteinte quand tous les processus du
          domaine atteignent le niveau de capacité requis.
\end{itemize}

\begin{boiteRemarque}{Et le modèle de COBIT~5~?}
COBIT~5 mesurait la capacité selon la norme ISO/IEC~15504, devenue
ISO/IEC~33000, dont les niveaux ont un sens assez différent. COBIT~2019
lui substitue un modèle inspiré du CMMI. Les organisations qui veulent
conserver l'approche ISO/IEC~33000 le peuvent~: COBIT~2019 précise que ses
pratiques tiennent lieu de « pratiques de base » et ses flux d'information
de « produits de travail ». Il faut seulement ne pas mélanger les deux
échelles dans un même rapport.
\end{boiteRemarque}

% ============================================================================
\section{Les Niveaux de Capacité}
\label{sec:c08-niveaux}
% ============================================================================

\subsection{Six niveaux}

\begin{boxdef}
\termecle{Capacité d'un processus}\index{capacité d'un processus}
(\textit{process capability})~: mesure de la manière dont un processus est
mis en œuvre et fonctionne. COBIT~2019 la mesure sur une échelle de~0
à~5.
\end{boxdef}

\begin{tableaucours}{cL{3.2cm}L{8.0cm}}{%
    \tableauHeader{Niv.} &
    \tableauHeader{Caractère} &
    \tableauHeader{Description}
}
0 & Absence & Aucune capacité de base~; approche incomplète~; les pratiques
ne sont pas, ou pas vraiment, mises en œuvre. \\
1 & Initial, intuitif & Le processus atteint plus ou moins sa finalité, par
un ensemble incomplet d'activités, peu organisé. \\
2 & Réalisé & Le processus atteint sa finalité par un ensemble d'activités
élémentaire, mais complet. \\
3 & Défini & Le processus atteint sa finalité de manière beaucoup plus
organisée, en s'appuyant sur des ressources de l'organisation~; il est
généralement bien défini. \\
4 & Mesuré & Le processus atteint sa finalité, il est bien défini, et sa
performance est mesurée de manière quantitative. \\
5 & Optimisé & Le processus est bien défini, sa performance est mesurée
pour l'améliorer, et l'amélioration continue est recherchée. \\
\end{tableaucours}
\vspace{-0.8em}
\begin{center}{\small\itshape Tableau~8.1. Les niveaux de capacité des
processus (d'après COBIT~2019, notre traduction)}\end{center}

\subsection{Des activités rattachées à des niveaux}

Le guide des objectifs rattache chaque activité à un niveau de capacité.
Dans la fiche DSS02, par exemple, les activités qui consistent à
enregistrer, classer et prioriser les demandes et les incidents
(pratique DSS02.02) relèvent du niveau~2~; celles qui définissent les
schémas de classification et les règles d'escalade (DSS02.01) relèvent du
niveau~3~; l'analyse des tendances et son utilisation pour l'amélioration
continue (DSS02.07) relèvent des niveaux~4 et~5~\cite{isaca2018cobitobjectifs}.

La progression est logique~: on commence par \emph{faire} (niveau~2),
puis on \emph{standardise} (niveau~3), puis on \emph{mesure} (niveau~4),
enfin on \emph{améliore} en continu (niveau~5).

\begin{boxalert}
\textbf{Violation typique.} Une organisation se déclare « au niveau~4 »
parce qu'elle produit un tableau de bord mensuel sur ses incidents. Mais
les incidents ne sont pas tous enregistrés, et chaque équipe les classe à
sa façon~: le tableau de bord mesure un processus qui n'est ni complet
(niveau~2) ni défini (niveau~3).

\smallskip
\textbf{Règle~:} un niveau ne s'acquiert que si les niveaux inférieurs
sont atteints. Mesurer un processus qui n'est pas défini produit des
chiffres, pas une capacité.
\end{boxalert}

% ============================================================================
\section{Capacité et Maturité}
\label{sec:c08-maturite}
% ============================================================================

\begin{boxdef}
\termecle{Maturité}\index{maturité} (\textit{maturity})~: mesure de
performance au niveau d'un domaine d'intérêt (un ensemble d'objectifs et
de leurs composantes). Un domaine d'intérêt atteint un niveau de maturité
donné lorsque \emph{tous} les processus qu'il contient atteignent ce
niveau de capacité.
\end{boxdef}

\begin{tableaucours}{cL{3.0cm}L{8.2cm}}{%
    \tableauHeader{Niv.} &
    \tableauHeader{Maturité} &
    \tableauHeader{Description}
}
0 & Incomplète & Le travail n'est pas forcément accompli. \\
1 & Initiale & Le travail est accompli, mais l'objectif du domaine n'est pas
pleinement atteint. \\
2 & Gérée & Planification et mesure de la performance existent, mais pas
encore de manière standardisée. \\
3 & Définie & Des standards communs à toute l'organisation guident le
travail. \\
4 & Quantitative & L'organisation est pilotée par les données, avec une
amélioration quantitative de la performance. \\
5 & Optimisée & L'organisation est tournée vers l'amélioration continue. \\
\end{tableaucours}
\vspace{-0.8em}
\begin{center}{\small\itshape Tableau~8.2. Les niveaux de maturité des
domaines d'intérêt (d'après COBIT~2019, notre traduction)}\end{center}

La règle d'agrégation est exigeante~: la maturité d'un domaine est celle
de son processus le plus faible. Une banque dont neuf processus de
sécurité sont au niveau~3 et le dixième au niveau~1 a une maturité de
niveau~1 dans le domaine de la sécurité. C'est voulu~: une chaîne a la
solidité de son maillon le plus faible.

\subsection{Mesurer les autres composantes}

COBIT~2019 reconnaît qu'il n'existe pas de méthode généralement admise
pour évaluer les composantes autres que les processus, mais il propose des
critères. Pour une structure organisationnelle, par exemple~: elle est
formellement créée~; son mandat est clair et documenté~; ses règles de
fonctionnement sont écrites~; elle se réunit comme prévu~; ses
procès-verbaux existent et sont utiles~; ses droits de décision sont
définis et respectés~; ses procédures d'escalade sont appliquées. Ces
critères rejoignent l'exigence de contrôle « un comité qui fonctionne » du
Chapitre~\ref{chap:gouv-entreprise}.

% ============================================================================
\section{Évaluer une Capacité}
\label{sec:c08-evaluer}
% ============================================================================

\subsection{L'échelle de notation}

Un niveau de capacité peut être atteint plus ou moins complètement.
COBIT~2019 propose, pour les évaluations orientées vers l'amélioration,
l'échelle suivante~\cite{isaca2018cobitintro}.

\begin{tableaucours}{L{3.6cm}L{3.6cm}L{5.0cm}}{%
    \tableauHeader{Note} &
    \tableauHeader{Taux de réalisation} &
    \tableauHeader{Lecture}
}
Complètement (\textit{Fully}, F) & Plus de 85~\% & Le niveau est acquis \\
Largement (\textit{Largely}, L) & De 50 à 85~\% & Acquis, avec des
faiblesses \\
Partiellement (\textit{Partially}, P) & De 15 à 50~\% & Non acquis \\
Pas (\textit{Not}, N) & Moins de 15~\% & Non acquis \\
\end{tableaucours}
\vspace{-0.8em}
\begin{center}{\small\itshape Tableau~8.3. L'échelle de notation des
niveaux de capacité (d'après COBIT~2019)}\end{center}

COBIT~2019 souligne que la note reste un jugement, qui doit être étayé par
l'examen des composantes~: activités réalisées, pièces produites,
structures en fonctionnement. Les méthodes formelles, qui débouchent sur
une certification indépendante, peuvent préférer une notation binaire
(atteint ou non).

\subsection{La règle d'attribution du niveau}

Le cours retient la règle usuelle des évaluations de processus, issue de
la famille ISO/IEC~33000~: \textbf{un processus atteint le niveau~$n$ si
ce niveau est noté au moins « largement » et si tous les niveaux
inférieurs sont notés « complètement »}. Un niveau noté « largement » est
donc acquis, mais il bloque l'accès au niveau suivant tant qu'il n'est pas
porté à « complètement ».

\subsection{Une évaluation pas à pas}

\begin{boxscenario}{Opération SANAGA. La capacité de BAI06}
L'auditrice interne, Mme~EKAMBI, évalue le processus BAI06 (gérer les
changements informatiques), dont la cible est le niveau~3. Elle examine,
pour chaque activité, les pièces de vingt changements tirés au hasard sur
six mois. \textit{(Le nombre d'activités par niveau est simplifié pour
l'exercice~; la fiche officielle BAI06 en donne la liste exacte.)}

\medskip
\begin{tableaucours}{cccccc}{%
    \tableauHeader{Niveau} &
    \tableauHeader{Activités} &
    \tableauHeader{Réalisées} &
    \tableauHeader{Taux} &
    \tableauHeader{Note} &
    \tableauHeader{Acquis~?}
}
2 & 10 & 6 & 60~\% & L & Oui \\
3 & 8 & 2 & 25~\% & P & Non \\
4 & 4 & 0 & 0~\% & N & Non \\
\end{tableaucours}

\medskip
\textbf{Constatation.} Six activités de niveau~2 sur dix sont réalisées~:
les demandes de changement sont enregistrées et évaluées, mais les
changements d'urgence ne font l'objet d'aucune procédure, et
l'autorisation n'est pas tracée pour les changements jugés mineurs. Deux
activités de niveau~3 sur huit sont réalisées.

\medskip
\textbf{Analyse} (confiance élevée). Capacité constatée~: \textbf{niveau~2,
largement atteint}. L'écart avec la cible est d'un niveau. Pour progresser,
il faut d'abord porter le niveau~2 à « complètement » (procédure
d'urgence, traçabilité de toutes les autorisations), puis réaliser les
activités de niveau~3.
\end{boxscenario}

\begin{boxprocedure}
\textbf{Une évaluation de capacité opposable.} Pour qu'une évaluation de
capacité puisse être présentée au conseil ou au régulateur, le dossier
doit contenir~:
\begin{itemize}
    \item l'objectif évalué, l'édition de COBIT et l'échelle de notation
          utilisées~;
    \item le périmètre (entités, systèmes, période) et la méthode
          d'échantillonnage~;
    \item pour chaque activité, la pièce examinée et le constat
          (réalisée ou non)~;
    \item le taux et la note de chaque niveau, et la règle d'attribution
          appliquée~;
    \item la capacité constatée, la capacité cible et l'écart~;
    \item la date, le nom de l'évaluateur et son indépendance par rapport
          au processus évalué.
\end{itemize}
\end{boxprocedure}

\begin{boxmetier}{Auditeur SI}
\textbf{Conduire une évaluation de capacité}
\begin{enumerate}
    \item Partir de la fiche officielle de l'objectif~: lister ses
          activités, niveau par niveau.
    \item Pour chaque activité, définir la pièce qui prouve sa réalisation
          (procès-verbal, ticket, journal, rapport).
    \item Échantillonner~: une activité réalisée une fois n'est pas une
          activité réalisée.
    \item Noter chaque niveau sur l'échelle F, L, P, N, puis appliquer la
          règle d'attribution.
    \item Restituer à l'audité avant de conclure~: il peut détenir une pièce
          que vous n'avez pas vue.
\end{enumerate}
\end{boxmetier}

% ============================================================================
\section{La Mise en Œuvre en Sept Phases}
\label{sec:c08-phases}
% ============================================================================

\subsection{Un programme, pas un projet}

Le guide de mise en œuvre de COBIT~2019 part d'un constat~: beaucoup de
démarches de gouvernance échouent parce qu'elles ne sont pas lancées et
pilotées comme de véritables \emph{programmes}, parrainés par la direction
générale, au périmètre maîtrisé et aux objectifs atteignables. Le but final
n'est pas de terminer un projet, mais d'installer une pratique durable~:
le programme se clôt quand l'amélioration de la gouvernance produit des
bénéfices mesurables et fait partie du fonctionnement normal.

\subsection{Sept phases, trois anneaux}

La démarche comporte sept phases, chacune formulée comme une question.
Chaque phase se lit sur trois anneaux~: le \emph{management du programme}
(anneau extérieur), l'\emph{accompagnement du changement} (anneau
intermédiaire) et le \emph{cycle d'amélioration continue} (anneau
intérieur)~\cite{isaca2018cobitintro}.

\begin{tableaucours}{L{3.3cm}L{2.9cm}L{2.9cm}L{3.1cm}}{%
    \tableauHeader{Phase} &
    \tableauHeader{Programme} &
    \tableauHeader{Changement} &
    \tableauHeader{Amélioration}
}
1. Quels sont les moteurs~? & Lancer le programme & Créer le désir de
changer & Reconnaître la nécessité d'agir \\
2. Où en sommes-nous~? & Définir problèmes et opportunités & Former
l'équipe de mise en œuvre & Évaluer l'état actuel \\
3. Où voulons-nous aller~? & Définir la feuille de route & Communiquer le
résultat visé & Définir l'état cible \\
4. Que faut-il faire~? & Planifier le programme & Identifier les acteurs &
Construire les améliorations \\
5. Comment y arriver~? & Exécuter le plan & Mettre en service et utiliser
& Mettre en œuvre les améliorations \\
6. Y sommes-nous arrivés~? & Réaliser les bénéfices & Ancrer les nouvelles
pratiques & Exploiter et mesurer \\
7. Comment garder l'élan~? & Examiner l'efficacité & Pérenniser &
Surveiller et évaluer \\
\end{tableaucours}
\vspace{-0.8em}
\begin{center}{\small\itshape Tableau~8.4. La feuille de route de mise en
œuvre de COBIT~2019 (d'après ISACA, notre traduction)}\end{center}

\subsection{Le lien avec la conception}

Les deux démarches s'emboîtent. La conception du système de gouvernance
(Chapitre~\ref{chap:cobit-conception}) nourrit les trois premières phases
de la mise en œuvre~: la compréhension du contexte alimente la phase~1,
la détermination et l'affinage du périmètre alimentent la phase~2, et la
conclusion de la conception fixe l'état cible de la phase~3. Les
évaluations de capacité, enfin, mesurent l'état actuel en phase~2 et les
progrès en phase~6.

% ============================================================================
\section{Construire une Feuille de Route}
\label{sec:c08-feuille}
% ============================================================================

\subsection{De l'écart au plan}

La phase~3 compare l'état cible et l'état actuel, puis cherche des
solutions. COBIT recommande de distinguer les \emph{gains rapides} des
chantiers de long terme, de donner la priorité aux actions faciles et à
fort bénéfice, et de découper les chantiers longs en étapes maîtrisables.

\begin{boxscenario}{Opération SANAGA. La feuille de route à dix-huit mois}
Mme~EKAMBI reprend son tableau sur ces bases. Il ne compte plus quarante
cases rouges, mais une ligne par objectif prioritaire. Voici les sept
objectifs du profil de conception (Chapitre~\ref{chap:cobit-conception})~;
EDM01, DSS02 et MEA03 suivent le même format.

\medskip
\begin{tableaucours}{L{1.5cm}cccL{6.2cm}}{%
    \tableauHeader{Objectif} &
    \tableauHeader{Constatée} &
    \tableauHeader{Cible} &
    \tableauHeader{Écart} &
    \tableauHeader{Première action}
}
DSS04 & 1 & 4 & 3 & Tester le plan de continuité du système central
(gain rapide~: 3~mois) \\
EDM03 & 1 & 4 & 3 & Faire approuver l'appétence au risque par le conseil \\
APO12 & 2 & 4 & 2 & Registre des risques SI relié au registre de la banque \\
BAI06 & 2 & 3 & 1 & Procédure des changements d'urgence (gain rapide) \\
APO10 & 1 & 3 & 2 & Renégocier le contrat de DIGICORE AFRIQUE \\
BAI07 & 1 & 3 & 2 & Recette et retour arrière obligatoires avant bascule \\
MEA02 & 2 & 3 & 1 & Suivi trimestriel des recommandations au comité
d'audit \\
\end{tableaucours}

\medskip
\textbf{Échelonnement.} Mois 1 à 3~: gains rapides (DSS04, BAI06, MEA02).
Mois 4 à 12~: chantiers de fond (EDM03, APO12, APO10, BAI07). Mois 13 à
18~: mesure des progrès (phase~6) et nouvelle évaluation de capacité.
\end{boxscenario}

\begin{boxmetier}{Responsable GRC}
\textbf{Piloter la feuille de route}
\begin{enumerate}
    \item Donner à chaque ligne un propriétaire, une échéance et un
          indicateur de progrès.
    \item Rendre compte chaque trimestre au comité stratégique des SI, et
          chaque semestre au conseil.
    \item Réévaluer la capacité des objectifs traités à mi-parcours, par une
          personne indépendante du processus.
    \item Mettre à jour la conception (Chapitre~\ref{chap:cobit-conception})
          si un facteur change~: c'est le principe du système dynamique.
\end{enumerate}
\end{boxmetier}

\begin{boiteRemarque}{Viser le niveau~5 partout~?}
Rien ne serait plus coûteux, ni plus inutile. Le niveau~5 suppose une
amélioration continue fondée sur des mesures quantitatives~: il se
justifie pour quelques processus réellement critiques, rarement plus. Pour
la plupart des organisations camerounaises, atteindre durablement le
niveau~3 sur les objectifs prioritaires serait déjà un progrès décisif~:
des processus définis, appliqués partout, qui ne dépendent plus de la
bonne volonté d'une personne. La gouvernance, ici encore, est affaire de
proportion.
\end{boiteRemarque}

% ============================================================================
\separateur
% ============================================================================

\begin{boxresume}
\textbf{Ce qu'il faut retenir de ce chapitre~:}
\begin{itemize}
    \item \textbf{Capacité et maturité}~: la capacité se mesure processus
          par processus (0 à 5)~; la maturité se mesure par domaine
          d'intérêt, et vaut celle du processus le plus faible.
    \item \textbf{Six niveaux}~: absence, initial, réalisé, défini, mesuré,
          optimisé. On fait, on standardise, on mesure, on améliore.
    \item \textbf{Des activités rattachées à des niveaux}~: c'est ce qui rend
          la mesure objective, à condition de l'étayer par des pièces.
    \item \textbf{Quatre notes}~: complètement (plus de 85~\%), largement
          (50 à 85~\%), partiellement (15 à 50~\%), pas (moins de 15~\%).
          Un niveau ne s'acquiert que si les niveaux inférieurs sont
          complètement atteints.
    \item \textbf{Sept phases, trois anneaux}~: programme, changement,
          amélioration continue. La gouvernance s'installe par un
          programme, pas par un projet.
    \item \textbf{La feuille de route} part des écarts entre capacité
          constatée et cible, commence par des gains rapides, et se pilote
          devant le comité stratégique.
\end{itemize}
\end{boxresume}

\bigskip

\begin{boxexo}[title={\ding{45}\quad Exercice~8.1. Noter un niveau \niveaubase}]
Pour chacun des cas suivants, donnez la note (F, L, P ou N) du niveau
concerné, puis la capacité du processus.
\begin{enumerate}[(a)]
    \item Niveau~2~: 9 activités sur 10 réalisées~; niveau~3~: 3 sur 6.
    \item Niveau~2~: 7 activités sur 10 réalisées~; niveau~3~: 6 sur 6.
    \item Niveau~2~: 1 activité sur 10 réalisée.
    \item Niveau~2~: 12 sur 12~; niveau~3~: 8 sur 9~; niveau~4~: 1 sur 5.
\end{enumerate}
\end{boxexo}

\begin{boxexo}[title={\ding{45}\quad Exercice~8.2. La maturité de la continuité \niveaumoyen}]
Le domaine d'intérêt « continuité » de KRIBILOG SA regroupe quatre
processus. Leurs capacités constatées sont~: DSS04 niveau~2, BAI04
niveau~3, DSS01 niveau~3, APO10 niveau~1.
\begin{enumerate}[(a)]
    \item Quelle est la maturité du domaine~? Justifiez.
    \item Le directeur général affirme que « la continuité est au niveau~3
          en moyenne ». Que lui répondez-vous~?
    \item Quelle action aurait le plus d'effet sur la maturité du domaine~?
\end{enumerate}
\end{boxexo}

\begin{boxexo}[title={\ding{45}\quad Exercice~8.3. La feuille de route de l'ANPP \niveauexpert}]
À partir des capacités cibles calculées à l'Exercice~7.2, et des capacités
constatées suivantes~: EDM02 niveau~0, APO05 niveau~1, APO10 niveau~1,
BAI11 niveau~2, DSS04 niveau~1.
\begin{enumerate}[(a)]
    \item Construisez le tableau des écarts.
    \item Identifiez deux gains rapides et deux chantiers de long terme.
    \item Organisez la feuille de route selon les sept phases, en précisant
          pour chaque phase une action sur chacun des trois anneaux.
    \item Indiquez comment la feuille de route s'inscrit dans le budget
          programme de l'agence (Chapitre~\ref{chap:alignement}).
\end{enumerate}
\end{boxexo}

\bigskip

\begin{boxinfo}
\textbf{Transition.} La Partie~\ref{part:cobit} a construit, avec COBIT~2019,
un système de gouvernance mesuré et une feuille de route. La
Partie~\ref{part:normatif} replace ce système dans son cadre normatif et
juridique. Le Chapitre~\ref{chap:normes} présente les normes
internationales de la gouvernance et des services (ISO/IEC~38500,
ISO/IEC~20000-1, ISO/IEC~27001, ISO~31000, COSO).
\end{boxinfo}
